\section{Model}\label{sec:model}

The primitives are a procurement auction for a single indivisible item, a set of potential bidders indexed by $i = 1, \ldots, N$, a private cost of supply $c_i$ drawn independently across bidders from a type-specific distribution $F_c^k$ where $k \in \{\text{SME}, \neg\}$ indexes bidder type, and an auction format. I take the auction format to be a Preg\~ao, which BEC implements as a descending-clock English-reverse auction: the clock opens at the reference price, firms can lower their offer at will, and the auction ends when the last firm willing to supply at the running clock price remains. Under the standard English-reverse interpretation \citep{haile2003}, bidders' weakly dominant strategy is to drop out when the clock reaches their cost, so losers' exit prices are point identifiers of their costs. The winner, by contrast, stays until the clock stops; her exit price is an upper bound on her cost ($c_{\text{win}} \le c_{(2)}$, the second-lowest exit).

\begin{assumption}[IPV, asymmetric across types]\label{a:ipv}
Costs $c_i$ are drawn independently across bidders. Conditional on type $k$, $c_i \sim F_c^k$ with density $f_c^k$ having full support on $[\underline{c}, \overline{c}] \subset \mathbb{R}_+$. The SME and non-SME distributions are not assumed to coincide.
\end{assumption}

\begin{assumption}[Auction-level multiplicative heterogeneity]\label{a:uh}
The observed normalized bid is $y_{it} = \log(b_{it} / p_t^{\mathrm{ref}}) = a_t + e_{it}$, where $a_t$ captures unobserved auction-level scale common to all bidders in auction $t$, and $e_{it}$ is bidder-specific and independent across bidders within $t$. $a_t$ and $e_{it}$ are mean-zero with finite variances $\sigma_a^{2,k}$ and $\sigma_e^{2,k}$ possibly varying by stratum.
\end{assumption}

Assumption~\ref{a:uh} is the Krasnokutskaya (2011) restriction that makes UH separable from the private value. It is an identifying restriction---$a_t$ and $e_{it}$ would not be separable without it---and is weaker than the full Kotlarski (1967) restriction, which would additionally require $a_t$ and $e_{it}$ to be symmetric within type. A Turnbull robustness check that does not impose Assumption~\ref{a:uh} but treats the winner as left-censored is reported in Section~\ref{sec:robustness}; the main estimates and the NPMLE agree within five percent.

\begin{assumption}[SME-only set-aside restricts the pool]\label{a:setaside}
Under the open regime, a non-empty subset of both types enters. Under the set-aside regime, only SMEs are admitted; the non-SME mass is truncated from the pool. Type-specific cost distributions $F_c^{\text{SME}}$ do not shift as a function of the regime change itself within a short window around the cutoff (cost-primitive stability).
\end{assumption}

Assumption~\ref{a:setaside} is what the policy change does in the model: it changes which types can enter, but not the cost distribution of a given type. This is the structural analogue of the identifying assumption in the reduced-form DiD---comparability between Pre and Post within type. Its empirical defense is the primitive-invariance KS test in Table~\ref{tab:v3_primitive_invariance}: the cost distribution of non-SMEs in pharma Preg\~ao is statistically invariant Pre-to-Post (D=0.032), validating that in the subsample where the assumption can be tested it holds. In non-pharma and pharma SME strata the invariance is noisier, and we rely on the cross-modality convergence (Convite GPV vs.\ Preg\~ao drop-out) reported in Table~\ref{tab:v3_cross_modality} as the main consistency check.

Equilibrium in the BNE with asymmetric IPV is characterized in \citet{maskinriley2000}. For the counterfactual simulations I exploit the fact that in a descending-clock auction with IPV, the winning price coincides in expectation with the second-lowest cost draw ($p = E[c_{(2)}]$) by revenue equivalence \citep{vickrey1961,krasnokutskaya2011}. This lets me simulate counterfactual prices by drawing from the estimated $F_c^k$, sorting, and reading off the second-lowest cost. The same strategy is used in the welfare analysis to compute production cost $c_{(1)}$ and allocative waste.

Entry is endogenous in the \citet{atheyseira2011} sense: I take the observed Pre-period counts $(N^{\text{SME}}_{\text{Pre}}, N^{\neg}_{\text{Pre}})$ as the status-quo arrival rates and the Post-period $(N^{\text{SME}}_{\text{Post}}, 0)$ as the equilibrium arrival under the restricted regime. This is a reduced-form treatment of entry---I do not estimate a structural fixed-cost-of-entry parameter---but the BNE counterfactual takes the equilibrium entry counts as data. Entry is modeled as Poisson with stratum-specific rate; sensitivity to alternative arrival assumptions is in the robustness battery.

With these pieces, the policy decomposition is well-defined. Let $p_{S_1}$ denote the simulated mean $c_{(2)}$ under the Pre-policy open regime, $p_{S_2}$ under SME-only with the Pre-period SME pool, and $p_{S_3}$ under SME-only with the Post-period SME pool. The intensive effect is $p_{S_2} - p_{S_1}$ (same pool size for SMEs, same absence of non-SMEs, $F_c$ fixed), and the entry effect is $p_{S_3} - p_{S_2}$ (pool expands to the observed post-policy level, $F_c^{\text{SME}}$ possibly shifted). The decomposition is the central quantitative object of the paper.
